\subsection{四元数代数的一般性质}

设 \(\alpha, \beta, \gamma\) 是 K 的元素，F 是型为 \((\alpha, \beta, \gamma)\) 的四元数代数。回忆（III, §2, No.~5, 第 445 页）\(^{(1)}\)：F 是一个结合含幺 K-代数，它在 K 上具有一组满足下列关系的基 \((1, i, j, k)\)。

\[
i ^ {2} = \alpha + \beta i, j ^ {2} = \gamma , i j = k, j i = \beta j - k.\tag{1}
\]

它是一个凯莱代数（III, §2, No.~4, 第 441 页，定义 1），其共轭满足

\[
\overline {{i}} = \beta - i, \quad \overline {{j}} = - j, \quad \overline {{k}} = - k.\tag{2}
\]

回忆 F 的凯莱迹与凯莱范数是从 F 到 K 的映射 \(T_{F}\) 与 \(N_{F}\)，它们由 \(\mathrm{T}_{\mathrm{F}}(q)=q+\overline{q}\) 与 \(\mathrm{N}_{\mathrm{F}}(q)=q\overline{q}\) 定义。

F 的以 \((1,i)\) 为基的线性子空间 E 是 F 的一个交换凯莱子代数；它是一个型为 \((\alpha,\beta)\) 的二次代数，且 F 可等同于由 \(\gamma\) 定义的 E 的凯莱扩张（III, §2, No.~5, 第 444 页）。对每个 \(z\in E\)，我们有 \(zj=j\overline{z}\)。F 的每个元素 q 都可唯一地写成 \(x+jy\)（其中 \(x,y\in E\)），且我们有

\[
\overline {{q}} = \overline {{x}} - j y, \quad \mathrm{T} _ {\mathrm{F}} (q) = x + \overline {{x}}, \quad \mathrm{N} _ {\mathrm{F}} (q) = x \overline {{x}} - \gamma y \overline {{y}}.\tag{3}
\]

命题 1. —— F 的元素 q 的特征多项式等于 \(\left(\mathrm{X}^{2}-\mathrm{T}_{\mathrm{F}}(q)\mathrm{X}+\mathrm{N}_{\mathrm{F}}(q)\right)^{2}\)。

由上所述，代数 \(\mathrm{F}\) 是一个自由右 \(\mathrm{E}\) -模，其基为 \((1,j)\)。因此，\(\mathrm{F}[\mathrm{X}]\) 是一个自由右 \(\mathrm{E}[\mathrm{X}]\) -模，其基为 \((1,j)\)。我们用 \(u\) 表示右 \(\mathrm{E}[\mathrm{X}]\) -模 \(\mathrm{F}[\mathrm{X}]\) 的、由 \(u(\mathrm{P}) = (\mathrm{X} - q)\mathrm{P}\)（对 \(\mathrm{P} \in \mathrm{F}[\mathrm{X}]\)）定义的自同构。\(q\) 的特征多项式是把 \(u\) 视为 \(\mathrm{K}[\mathrm{X}]\) -模 \(\mathrm{F}[\mathrm{X}]\) 的自同构时的行列式。由 III, §9, No.~4, 第 546 页，命题 6，它等于 \(\mathrm{N}(\det u)\)，其中 \(\mathrm{N}\) 表示从 \(\mathrm{E}[\mathrm{X}]\) 到 \(\mathrm{K}[\mathrm{X}]\) 的范数。把 \(q\) 写成 \(x + jy\)（其中 \(x, y \in \mathrm{E}\)）。\(u\) 关于基 \((1,j)\) 的矩阵是 \(\begin{pmatrix} \mathrm{X} - x & -\gamma \overline{y} \\ -y & \mathrm{X} - \overline{x} \end{pmatrix}\)；它的行列式等于 \(\mathrm{D} = (\mathrm{X} - x)(\mathrm{X} - \overline{x}) - \gamma y\overline{y} = \mathrm{X}^2 - \mathrm{T}_{\mathrm{F}}(q)\mathrm{X} + \mathrm{N}_{\mathrm{F}}(q)\)（参看公式 (3)）。由于 \(\mathrm{D}\) 属于 \(\mathrm{K}[\mathrm{X}]\)，我们有 \(\mathrm{N}(\mathrm{D}) = \mathrm{D}^2\)；命题 1 得证。

注. —— 1) 设 K 的特征异于 2，并令 \(i' = 2i - \beta\)。则 \((1, i')\) 是 E 在 K 上的一组基，且我们有 \(i'^2 = 4\alpha + \beta^2\)。由此推出，E 同构于型为 \((4\alpha + \beta^2, 0)\) 的二次代数，F 同构于型为 \((4\alpha + \beta^2, 0, \gamma)\) 的四元数代数。

\begin{enumerate}
\def\labelenumi{\arabic{enumi})}
\setcounter{enumi}{1}
\item
  型为 \((\alpha,0,\gamma)\) 的四元数代数同构于型为 \((\gamma,0,\alpha)\) 的四元数代数（III, §2, No.~5, 第 445 页）。它也同构于型为 \((\alpha a^{2},0,\gamma c^{2})\) 的四元数代数（对 K 的非零元素的每一对 \((a,c)\)）。
\item
  设 q 是 F 的一个元素。则 q 是幂零的，当且仅当它的特征多项式等于 \(X^{4}\)，即 \(\mathrm{T}_{\mathrm{F}}(q)=\mathrm{N}_{\mathrm{F}}(q)=0\)；这时我们有 \(q^{2}=0\)。
\end{enumerate}

例. —— 矩阵代数 \(\mathbf{M}_{2}(\mathrm{K})\) 同构于型为 \((0,1,1)\) 的四元数代数。事实上，考虑二次代数 \(E = K \times K\)（其型为 \((0,1)\)）与四元数代数 \(F = E + Ej\)（它是由元素 \(\gamma = 1\) 定义的 E 的凯莱扩张）。映射 \((a,b) \mapsto \begin{pmatrix} a & 0 \\ 0 & b \end{pmatrix}\) 是从 E 到 \(\mathbf{M}_{2}(\mathrm{K})\) 的一个代数同态。由于对 K 中的 a, b，我们有

\[
\left( \begin{array}{c c} 0 & 1 \\ 1 & 0 \end{array} \right) \left( \begin{array}{c c} 0 & 1 \\ 1 & 0 \end{array} \right) = \left( \begin{array}{c c} 1 & 0 \\ 0 & 1 \end{array} \right), \quad \left( \begin{array}{c c} 0 & 1 \\ 1 & 0 \end{array} \right) \left( \begin{array}{c c} a & 0 \\ 0 & b \end{array} \right) = \left( \begin{array}{c c} b & 0 \\ 0 & a \end{array} \right) \left( \begin{array}{c c} 0 & 1 \\ 1 & 0 \end{array} \right),
\]

这个同态扩张为一个代数同态 \(\theta : \mathrm{F} \longrightarrow \mathbf{M}_2(\mathrm{K})\)，其定义为

\[
\theta \bigl ((a, b) + (c, d) j \bigr) = \left( \begin{array}{c c} a & c \\ d & b \end{array} \right)  .
\]

这个同态是双射。当 \(\mathbf{K}\) 的特征异于 2 时，代数 \(\mathbf{M}_2(\mathbf{K})\) 也同构于型为 (1,0,1) 的四元数代数（注 1）。

